\begin{ouverture}[probabilites]{11}{Probabilités et expériences à deux étapes}{Partie 3 · Données et hasard}
{Issues · Événements · Tableaux · Arbres · Fréquences}
\objectif{Calculer dans une situation d’équiprobabilité.}
\objectif{Énumérer les issues d’une expérience à deux étapes.}
\objectif{Distinguer probabilité et fréquence observée.}
\prerequis{Fractions et proportionnalité.}
\end{ouverture}
\pagestyle{otMain}\essentiel\fontsize{11}{15}\selectfont
\essSection{1}{Modéliser le hasard}
\begin{definition}Une issue est un résultat possible. Un événement est un ensemble d’issues. Sa probabilité est comprise entre 0 et 1. Un événement impossible a pour probabilité 0 ; l’événement certain a pour probabilité 1.\end{definition}
\begin{propriete}Si toutes les issues sont équiprobables, la probabilité d’un événement vaut le nombre d’issues favorables divisé par le nombre total d’issues. Justifier cette équiprobabilité : un dé équilibré, des boules indiscernables au toucher, un choix uniforme.\end{propriete}
\begin{propriete}L’événement contraire de A contient les issues qui ne réalisent pas A. Sa probabilité est $1-P(A)$. Si deux événements ne peuvent se produire ensemble, la probabilité de leur réunion est la somme de leurs probabilités.\end{propriete}
\begin{attention}Compter des noms de couleurs n’est pas compter des issues équiprobables si les couleurs n’ont pas les mêmes effectifs. « Au moins un » inclut un ou plusieurs.\end{attention}
\clearpage
\essSection{2}{Deux étapes et simulations}
\begin{savoirfaire}[Énumérer]Un tableau à double entrée ou un arbre permet de ne rien oublier. Une feuille de l’arbre représente une succession ordonnée d’issues. Les branches issues d’un même nœud décrivent toutes les possibilités à cette étape.\end{savoirfaire}
\begin{exemple}Deux pièces équilibrées lancées indépendamment donnent quatre issues équiprobables : PP, PF, FP, FF. L’événement « exactement une face F » contient PF et FP, donc a pour probabilité $2/4=1/2$.\end{exemple}
\begin{savoirfaire}[Utiliser des poids]Pour les expériences décrites ici, multiplier les probabilités successives le long d’un chemin, puis additionner les chemins favorables. Sans remise, les nombres de boules et le total changent après le premier tirage. L’arbre explicite ce changement, sans formalisme de probabilité conditionnelle.\end{savoirfaire}
\begin{propriete}La fréquence observée est le nombre de succès divisé par le nombre d’essais. Pour beaucoup de répétitions dans les mêmes conditions, elle tend généralement à se stabiliser près de la probabilité. Cela ne prédit ni le prochain résultat ni une égalité exacte pour un nombre fini d’essais.\end{propriete}
\clearpage\rubrique{Découverte}
\begin{blocLivre}{ex-11-1}\begin{exercice}[1]{Un dé}
On lance un dé équilibré à six faces numérotées de 1 à 6. Quelle est la probabilité d’obtenir un nombre pair ?
\end{exercice}\end{blocLivre}
\begin{blocLivre}{ex-11-2}\begin{exercice}[1]{Des boules colorées}
Un sac contient 5 boules rouges, 3 bleues et 2 vertes, indiscernables au toucher. On tire une boule au hasard. Donner les probabilités de rouge et de non-rouge.
\end{exercice}\end{blocLivre}
\begin{blocLivre}{ex-11-3}\begin{exercice}[1]{Impossible ou certain}
Sur un dé classique, donner les probabilités d’obtenir 9 et d’obtenir moins de 7.
\end{exercice}\end{blocLivre}
\begin{blocLivre}{ex-11-4}\begin{exercice}[1]{Une fréquence expérimentale}
On obtient 46 fois face au cours de 100 lancers d’une pièce. Donner la fréquence de face. Est-ce une preuve que la pièce est déséquilibrée ?
\end{exercice}\end{blocLivre}
\begin{blocLivre}{ex-11-5}\begin{exercice}[1]{Une roue à secteurs égaux}
Une roue a huit secteurs de même angle numérotés de 1 à 8 ; tous ont la même chance d’arrêt. Quelle est la probabilité d’un multiple de 3 ?
\end{exercice}\end{blocLivre}

\clearpage\rubrique{Entraînement}
\begin{blocLivre}{ex-11-6}\begin{exercice}[2]{Deux pièces}
Deux pièces équilibrées sont lancées indépendamment. Compléter la liste des quatre issues et calculer la probabilité d’exactement une face F.
\label{fig-ex-11-6}\figureGeo{proba-pieces}
\end{exercice}\end{blocLivre}
\begin{blocLivre}{ex-11-7}\begin{exercice}[2]{La somme de deux dés}
Deux dés équilibrés sont lancés indépendamment. Utiliser le tableau des sommes pour calculer la probabilité d’une somme égale à 7.
\label{fig-ex-11-7}\figureGeo{proba-des}
\end{exercice}\end{blocLivre}
\begin{blocLivre}{ex-11-8}\begin{exercice}[2]{Avec remise}
Un sac contient 3 boules rouges et 2 bleues. On tire, on remet la boule, puis on tire de nouveau au hasard. Calculer la probabilité de deux rouges.
\label{fig-ex-11-8}\figureGeo{proba-remise}
\end{exercice}\end{blocLivre}
\begin{blocLivre}{ex-11-9}\begin{exercice}[2]{Sans remise}
Dans le même sac de 3 rouges et 2 bleues, on tire deux boules successivement sans remise. Calculer la probabilité de deux rouges.
\label{fig-ex-11-9}\figureGeo{proba-sans-remise}
\end{exercice}\end{blocLivre}
\begin{blocLivre}{ex-11-10}\begin{exercice}[2]{Au moins un six}
Deux dés équilibrés sont lancés indépendamment. Calculer la probabilité d’obtenir au moins un six.
\end{exercice}\end{blocLivre}

\clearpage\rubrique{Problèmes}
\begin{blocLivre}{ex-11-11}\begin{exercice}[3]{Deux couleurs différentes}
Reprendre le tirage sans remise de l’exercice 9. Calculer la probabilité d’obtenir deux couleurs différentes.
\label{fig-ex-11-11}\figureGeo{proba-sans-remise}
\end{exercice}\end{blocLivre}
\begin{blocLivre}{ex-11-12}\begin{exercice}[3]{Un jeu à points}
On paie 2 points pour lancer un dé équilibré. Le 6 rapporte 8 points, le 5 rapporte 4 points, les autres 0. Sur un grand nombre de parties, le gain brut moyen théorique est la somme des gains multipliés par leurs probabilités. Le jeu équilibre-t-il mise et gain moyen ?
\end{exercice}\end{blocLivre}
\begin{blocLivre}{ex-11-13}\begin{exercice}[3]{Un effectif inconnu}
Un sac contient 6 rouges et $b$ bleues seulement. La probabilité de tirer rouge au hasard est $2/5$. Trouver $b$.
\end{exercice}\end{blocLivre}
\begin{blocLivre}{ex-11-14}\begin{exercice}[3]{Un code à deux chiffres}
Un générateur choisit indépendamment et uniformément deux chiffres parmi 0, 1, 2, 3, avec répétition possible. Un code peut commencer par 0. Quelle est la probabilité que ses deux chiffres soient distincts ?
\end{exercice}\end{blocLivre}
\begin{blocLivre}{ex-11-15}\begin{exercice}[3]{Estimer sans promettre}
Une simulation d’un événement de probabilité $1/3$ donne 327 succès en 1000 essais, puis 3354 en 10000 essais. Comparer les écarts entre fréquences et probabilité. Peut-on garantir que tout essai supplémentaire rapproche la fréquence de $1/3$ ?
\end{exercice}\end{blocLivre}

\clearpage
